\section[Linguistic measurement and executable interventions]{Linguistic measurement and executable interventions in banking}
\label{sec:exec28-lm}

A language-assisted banking decision begins with an intervention that a bank can actually implement. A complaint team may decide whether to rewrite an acknowledgement, request clarification or escalate a case. An analyst may call the treatment politeness, yet this word can denote the writer's intention, the customer's perception, a classifier score or a particular editing procedure. These objects need not coincide. Their separation is necessary before an agent translates the business question into a graph, proposes an adjustment set or calls an estimator. A well-formed causal query about the wrong object remains the wrong business answer.

\subsection{Intent, perception, labels and actions}

The inspected preprint of \citet{exec28_lm_pryzant} distinguishes writer intent $T$, reader perception $\widetilde T$, a proxy label $\widehat T$ and text $W$. Its first identification result assumes that $W$ blocks the relevant backdoor paths and that intent agrees with perception, $T=\widetilde T$. These are substantive assumptions. A bank may intend reassurance while a distressed customer perceives evasion; a classifier may score deferential wording highly even though the letter makes no useful commitment. Neither a fluent explanation nor high predictive accuracy establishes agreement between these variables.

An operational treatment should identify an editing rule, the eligible message population, the permissible information and the distribution of text versions produced by the rule. For example, a rule might replace a discretionary closing paragraph with a fixed, reviewed clarification of the next servicing step. It must specify whether facts, deadlines and promised remedies can change. If these attributes vary together with tone, the resulting policy effect concerns the entire rewriting procedure. Calling it an isolated effect of politeness requires further restrictions; the proposed product should not silently impose them.

The same distinction guides diagram construction. Pre-existing case severity can influence both message selection and subsequent resolution. A customer response following the edited message is downstream of that action, even if an agent can later extract it into a numerical feature. Putting that response into a baseline adjustment set changes the question and can invalidate identification. The assistant should retain a variable's event time, observation time, decision availability and semantic definition alongside its name. It can suggest candidate arrows and competing explanations, while an explicit causal model states which claims are assumed and what evidence could contradict them.

The source's linguistic formulation provides a conceptual foundation rather than a certificate for an arbitrary embedding. Appendix A discusses the difficulty that a representation can perfectly predict a proxy treatment, leaving no overlap for a within-representation comparison. Its consistency discussion assumes an outcome-sufficient component, appropriate support and consistent regression. The preprint's theorem and surrounding notation alternate perception and proxy symbols; we retain its stated conditions without presenting a silently repaired theorem. A compression method must justify both preservation of adjustment information and usable treatment support. Reducing dimensionality alone proves neither.~\citep{exec28_lm_pryzant}

\subsection{Why classifier accuracy is insufficient}

Under the source's assumptions, define the conditional error probabilities
\begin{align}
 \epsilon_0(W)&=\Pr(\widetilde T=0\mid\widehat T=1,W),\\
 \epsilon_1(W)&=\Pr(\widetilde T=1\mid\widehat T=0,W).
\end{align}
Writing $\delta(W)$ for the conditional perception contrast, its proxy expression has the form
\begin{equation}
 \psi^{\mathrm{proxy}}=E\!\left[\{1-\epsilon_0(W)-\epsilon_1(W)\}\delta(W)\right].
 \label{eq:exec28-lm-proxy}
\end{equation}
The accompanying discussion explicitly qualifies directional attenuation by a better-than-chance classifier and a treatment-effect sign that is homogeneous across possible texts. This qualification must survive translation into a product specification. An overall accuracy number does not ensure that the factor multiplying the contrast is well behaved in each important stratum, or that the underlying effects have a common sign.~\citep{exec28_lm_pryzant}

A separate, fully stipulated measurement example makes the consequence concrete. Let $S\in\{0,1\}$ denote two equally prevalent servicing populations. Within each population, let a latent perception $R$ be assigned with probability one half. Define potential monetary outcomes by
\begin{equation}
 Y(r)=100+\delta_S r,\qquad \delta_0=20,\quad\delta_1=-10,
 \label{eq:exec28-lm-witness}
\end{equation}
in USD per account. Generate an observed proxy $Z$ by independently flipping $R$, conditional on $S$, with probabilities $q_0=9/20$ and $q_1=0$. These are authored assumptions, not estimates from complaints. The proxy is stochastic given the adjustment stratum; this is not the source's deterministic full-text proxy setting or an empirical replication of TextCause.

The true average effect is $(20-10)/2=5$ dollars. Symmetry gives $\Pr(Z=1\mid S)=1/2$. Bayes' rule yields $E[R\mid Z=1,S=s]=1-q_s$ and $E[R\mid Z=0,S=s]=q_s$. Consequently,
\begin{equation}
 E[Y\mid Z=1,S=s]-E[Y\mid Z=0,S=s]=(1-2q_s)\delta_s.
\end{equation}
The two proxy contrasts are $2$ and $-10$ dollars, averaging $-4$ dollars. Overall classification accuracy is
\begin{equation}
 \tfrac12(1-q_0)+\tfrac12(1-q_1)=\frac{31}{40}=77.5\%.
\end{equation}
Exact enumeration of all eight joint cells independently reproduces these conditional means and the algebraic calculation. For a stipulated portfolio of 10,000 accounts, the true gross contrast is USD~50,000 while the proxy comparison suggests minus USD~40,000, before implementation costs.~\citep{exec28_lm_reference}

This example violates the homogeneous-sign condition and therefore does not refute the source's conditional result. A high-accuracy semantic label may still give the wrong direction for a heterogeneous business effect. Evaluation should examine conditional errors, effect heterogeneity, population weighting and support rather than promote a single accuracy score into causal assurance. If perception cannot be measured reliably, the product can report sensitivity ranges or change the intervention to an observable, randomized editing policy. It should not invent a validated perception variable.

\subsection{Public complaints as a prototype resource}

The source supplies two distinct kinds of evaluation. Its Amazon setting constructs semi-synthetic outcomes, permitting comparison within the stipulated outcome model. Its CFPB application uses observational complaints and a fast-response indicator defined by a response in fewer than 15 days. It matches complaints with opposite outcomes, selects 4,000 similar pairs and retains the upper and lower quartiles of a politeness score, producing 4,000 selected complaints. These choices define a selected study population, not the entire customer population.~\citep{exec28_lm_pryzant}

The reported effect estimates should not be relabelled as known causal truth. Matching on outcomes and restricting to score extremes require an explicit account of selection and of the population to which a result applies. An agent-generated summary must preserve those choices rather than quote a favourable estimate as a verified benefit of adopting a rewriting service. Generalization to unselected complaints, contemporary servicing and repeated interactions requires additional justification.

For a proposal-stage prototype, public or generated data can establish useful facts at much lower cost than a bank experiment. A frozen public-data snapshot can exercise event-time parsing, missing-text handling, semantic label provenance and the distinction between response time and substantive resolution. A generated outcome construction can test whether an estimator recovers an effect under an explicit assignment and measurement law. An observed complaint archive provides realistic input structure, while a stipulated generator supplies known potential outcomes only within its model. Neither supplies otherwise unobserved counterfactual customer outcomes.

A prospective comparison should freeze the unit, sample selection, source release, permitted fields, intervention rule, estimator and all attempted requests before evaluating the candidate. Customer, case and temporal overlap should be checked when defining training and evaluation partitions. Prompts, retries, manual corrections, failed parses and excluded records contribute to both the denominator and the cost. These are proposed design requirements, not claims that an admitted runtime already satisfies them. A prototype can demonstrate the merit of a later, narrower bank trial without claiming authentic bank validation.

\subsection{From a one-message question to a longitudinal policy}

Repeated servicing introduces feedback. Earlier messages affect later customer responses, which influence subsequent messages and eventual outcomes. The retained CausalCollab work studies sequences of human edits and language-model replies and an incremental stylistic target. Its intervention definition, representations and generated evaluation design should be compared with an executable bank editing policy, rather than treated as interchangeable with a single-stage perception contrast.~\citep{exec13_census_collaboration}

For a separately specified finite-horizon policy, let $H_t$ contain information available immediately before decision $t$, $A_t$ be the editing or escalation action and $g_t(a\mid H_t)$ its prescribed probability. Let $L_{t+1}$ contain the next observations. With a terminal monetary outcome $Y$, the target is $V(g)=E[Y^g]$. Under consistency, sequential exchangeability, positivity for actions used by $g$, and a correctly stated observation process, a finite discrete longitudinal g-formula is
\begin{equation}
 V(g)=\sum_{\bar l,\bar a}m(h_K,a_K)\,p(l_0)
 \prod_{t=0}^{K}g_t(a_t\mid h_t)
 \prod_{t=0}^{K-1}p(l_{t+1}\mid h_t,a_t),
 \label{eq:exec28-lm-policy}
\end{equation}
where $m(h_K,a_K)=E[Y\mid H_K=h_K,A_K=a_K]$ and histories are chronological. Continuous variables require integration. This applies the longitudinal framework; it is not a new identification theorem for arbitrary language.~\citep{hernan2020}

The formula exposes the information the assistant must help specify. The action must denote a procedure, not a classifier label that retrospectively describes a message. Each history must contain the confounders claimed sufficient at that time; unobserved distress or operator judgement cannot be removed merely by generating a diagram. A deterministic action outside historical support invalidates positivity. Representation compression and estimated transition models introduce further approximation and learning questions. The single-stage measurement witness supplies no theorem resolving these longitudinal issues.

The graph-building assistant can produce a reviewable intervention contract: decision times, observable actions, eligible histories, competing graphs, proposed adjustment information, support diagnostics and unresolved assumptions. A reviewer should be able to trace a request to reduce avoidable repeat contact to a defined outcome and a contrast between implementable policies. An unsupported direction, ambiguous treatment or unavailable history should remain visible. This is more informative for prototype investment than a point estimate whose interpretation changed during translation.

\subsection{Adoption choices at the component boundary}

Scientific relevance and executable reuse answer different questions. Three pinned implementations illustrate why component-specific decisions are necessary. The inspected LLM-CD tree contains 2,899 entries and 2,745 files. Its nested BI-CaMML notice states redistribution conditions for that component, while the complete pinned tree has no root notice. Requirements specify Gurobi and pygobnilp versions; compiled artifacts also require correspondence to a reproducible source build. A nested notice and requirements file do not establish permission or a complete installed environment for the whole pipeline. Direct reuse remains deferred pending authoritative component terms, solver entitlement and artifact-to-source provenance.~\citep{exec28_lm_artifacts}

The inspected LLM-and-SCD snapshot has 15 entries and 12 files, without a root licence, README or environment manifest in that tree. Notebook and data paths are discovery leads, but do not establish execution order, full service identities or reproducibility against the paper. The proposal can retain the scientific approach as a comparator while deferring direct reuse until permissions, environment and source-to-result correspondence are available. A separately specified implementation would need its own scope, evaluation and source review.~\citep{exec28_lm_artifacts}

LCM presents a different boundary. Its root code notice is Apache 2.0, while four named pretrained-model links and three experiment-folder links point to external artifacts. Exact model bytes, hashes, provenance and applicable terms have not been established by the inspected code notice. Several dependency versions remain open, and vendored components need their own inventory. The inspected loader sets \texttt{weights\_only=False} and can return \texttt{None} after missing-file or exceptional loads; a caller must distinguish those outcomes from a valid model or a scientifically inconclusive result. An in-tree baseline checkpoint is not a substitute for the advertised LCM weights. These are static source observations, not executed failure rates or legal clearance.~\citep{exec28_lm_artifacts}

For CLI and MCP integration, the proposed request should carry the target and intervention contract, input schema and clock, selected model identity and artifact manifest. The response should separate an unsupported scientific target, an unavailable artifact, an execution failure, an inconclusive analysis and a successful result with declared assumptions. A fallback may be proposed as an explicit new method choice; silently replacing the intended model changes the comparison. A public fixture should test these branches and record all attempted requests and costs. Where external artifacts remain unavailable, independent work on an authorized estimator, public-data design or semantic contract can proceed. The proposed product must preserve the distinction between a meaningful scientific question, a justified answer and an executable method.
